\section{Erd\H{o}s Problem 231: finite words without abelian squares}

\subsection{1) Statement and scope}
Must a word of length $2^k-1$ on $k$ letters contain two adjacent
nonempty blocks with the same multiset of letters? The literal finite
statement has a negative answer for every $k\geq1$.
The construction below settles precisely this length formulation.
The related existence of an infinite abelian-square-free word on four
letters is a stronger theorem, due to Ker\"anen, and is not proved here.

\subsection{2) Recursive counterexample}
Let $W_1=1$ and, using a new letter at each stage, define
\[
 W_k=W_{k-1}\,k\,W_{k-1}\qquad(k\geq2).
\]
Induction gives $|W_k|=2^k-1$, and $W_k$ uses exactly the letters
$1,\ldots,k$. The first examples are
\[
 1,\qquad121,\qquad1213121,\qquad121312141213121.
\]
We prove by induction that $W_k$ contains no abelian square.
The assertion is immediate for $k=1$. Suppose that two adjacent
blocks $X,Y$ in $W_k$ have equal multisets. If their union contains
the central letter $k$, that letter occurs exactly once in $XY$.
It therefore cannot occur equally often in $X$ and $Y$, a contradiction.
If their union avoids $k$, the contiguous block $XY$ lies wholly
inside one of the two copies of $W_{k-1}$, contradicting induction.
This proves the claim.

\subsection{3) Interpretation and attribution}
The page's historical discussion suggests that a power-of-two length
may have been intended in an older formulation. Our word has length
$2^k-1$, not $2^k$; those questions must be distinguished.
Also, taking a limit in this recursion introduces infinitely many
different letters. It does not establish the infinite four-letter
result. This is an elementary reconstruction addressing the displayed
statement, with no claim of a new result or a new formal proof.

Source: T.~F.~Bloom, Erd\H{o}s Problem 231,
\url{https://www.erdosproblems.com/231}, accessed 23 September 2026.
The page attributes the infinite four-letter result to
V.~Ker\"anen (1992).

\subsection{4) Completion estimate}
The stated finite-length question is completely answered by the
explicit counterexample and induction.
\textbf{COMPLETION: 100\%}
